\begin{proof}

\begin{enumerate}[label={$\langle1\rangle$\arabic*.},leftmargin=*]
\item We retain the smooth curve $A$, its degree-two horizontal
cover $H$, and its degree-one vertical cover $R$. All other
normalized conductor components map to $B$. Let $r_B$ be
the generic normalization branch number along $B$, and
$c_B$ the number of global conductor components over it.
Their finite degrees sum to $r_B$. We exclude $r_B=1,2$.
The normalized conductor has $2+c_B$ components and a tree
graph, with total identification contribution $1+c_B$.
Distinct points on one normalized component cannot be
identified, since that would create a graph loop.

\paragraph{A length-two intersection calculation.}
Let $F,G$ be two local hypersurface branches whose
scheme-theoretic intersection is supported on the smooth
curve $A$ and has generic length $m\in\{1,2\}$ along it.
The ring
\[
                 Q=\mathbb C\{x,y,z\}/(F,G)
\]
is a one-dimensional complete intersection. Its only
associated prime is the ideal of $A$. Choose $z$ as a
coordinate on $A$. Then $Q$ is finite and torsion-free
over $\mathbb C\{z\}$, hence free of rank $m$. The
reduction map to $\mathbb C\{z\}$ splits through the
chosen coordinate.

\item For $m=1$, the ring is $\mathbb C\{z\}$, and the linear
parts of $F,G$ have rank two. Both branches are smooth.
For $m=2$, the nilradical is a free rank-one module over
$\mathbb C\{z\}$. Its square vanishes generically and
therefore vanishes, by torsion-freeness. A generator
$\epsilon$ gives the exact analytic algebra isomorphism
\[
                    Q\simeq\mathbb C\{z,\epsilon\}/(\epsilon^2).
                                                            \tag{4.56a}
\]
Thus $Q$ has embedding dimension two, so one of $dF,dG$
is nonzero. If $G$ has generic normal order two along $A$,
it belongs to the ordinary square of its smooth ideal and
$dG$ vanishes on all of $A$. Then $dF\ne0$ and $F$ is
smooth. We use this calculation on the original intersection
scheme, retaining its multiplicity.

\paragraph{Two distinct image junctions.}
At a two-branch image point, let $r_A,r'_B$ count the
normalized conductor points over its two branches, and
let $n$ count normalization points of $S$. Pairwise
intersection of the hypersurface branches gives a common
element in their supports in $\{A,B\}$. Hence
$n\le\max(r_A,r'_B)$, and the graph contribution is at
least $\min(r_A,r'_B)$. Here $r_A=2$ or three, from the
one or two $H$ points and the unique $R$ point.

\item If $c_B=2$, both degrees are one, and $r'_B=2$ at either
junction. The two contributions total at least four,
exceeding the budget three. If $c_B=1$ and its degree is
two, an unramified junction costs at least two and the
other at least one, exceeding the budget two. Thus one
component $V$ remains, with exactly one normalized point
over each junction.

\item If $B$ belongs to every branch support at a junction,
there is only one normalization point there and its
contribution is $r_A\ge2$. The other junction makes this
impossible. Therefore $A$ belongs to every support. The
budget forces contribution exactly one at each junction:
every normalization point has one $A$ branch, and the
sole $V$ point is paired with one of them. Connectedness
forces the graph path $H$--$V$--$R$, with its edges at
the two different junctions.

\item At the $R$--$V$ junction, let $F$ be the corresponding
hypersurface branch. The other branches come from $H$
and carry only $A$. If its two preimages are distinct,
intersection with either such branch has support $A$ and
generic length one, because the full generic $A$ germ
is an ordinary triple point. The preceding calculation
makes $F$ smooth. If the two $H$ preimages coalesce,
their one hypersurface branch $G$ has generic normal
order two along $A$, and the intersection with $F$ has
length two. Formula (4.56a) again makes $F$ smooth.

\item This is impossible. The branch $F$ carries all the nearby
generic $B$ preimages: either the single branch of a
singular unibranch transverse germ, or both branches
when $V\to B^\nu$ has degree two and ramifies at this
point. In either case $F$ is intrinsically nonnormal
along the nearby $B$ germ. It cannot be smooth in a
neighbourhood of the junction.

\paragraph{A three-branch image point with one $B$ component.}
The image branches are $A,B_+,B_-$, since $A$ is smooth.
If $c_B=1$, the normalized points on $V$ over the two
$B$ signs are distinct and cannot be identified in $N$.
No normalization point can therefore carry both signs.
Pairwise hypersurface intersections force every point
carrying either sign to carry $A$ as well.

\item The graph budget two now forces exactly one point over
each sign, each paired with exactly one $A$ branch, and
no other identification. The pairs are $H$--$V$ and
$R$--$V$. An extra $H$ preimage is $A$-only. At the
$R$--$V$ point, the other $H$ branch or branches intersect
its hypersurface branch only along $A$, because their
$B$ supports have the opposite sign or are empty. The
same length-one or length-two calculation makes this
branch smooth, contradicting its intrinsic nonnormality
along the $B$ germ whose generic preimages it contains.

\paragraph{A three-branch point with two $B$ components.}
Write the two degree-one components as $V,W$. If $h=1$
or two is the number of $H$ points over the $A$ value,
the total normalized branch count is
\[
                         T=(h+1)+2+2=h+5.
\]
The graph budget three gives $n\ge h+2$ normalization
points. No singleton image support is possible: a
singleton $\{A\}$ forces $n\le h+1$, and a singleton
$\{B_+\}$ or $\{B_-\}$ forces $n\le2$. Hence every
point has at least two image labels and $n\le T/2$.
This excludes $h=2$. For $h=1$ there are precisely
three points, each with exactly two conductor branches
of different image labels. Up to signs and interchange
of $V,W$, their supports are
\[
 \begin{array}{c|c}
 \text{hypersurface branch}&\text{conductor branches}\\ \hline
 F_R&R,\ B_-\text{ on }V\\
 F_H&H,\ B_+\text{ on }W\\
 F_0&B_+\text{ on }V,\ B_-\text{ on }W.
 \end{array}
                                                            \tag{4.56b}
\]
The tree is the path $H$--$W$--$V$--$R$, and $H\to A$
is ramified with local degree two at its point.

\item The intersection of $F_R$ and $F_H$ is supported on $A$
and has generic length two. The equation of $F_H$ has
normal order two along $A$, so (4.56a) makes $F_R$
smooth. Choose its equation as $X$, a coordinate $Y$
on that surface cutting $A$, and a coordinate $z$ on
$A$. Preparation and depression in $Y$, preserving the
ideal $(X,Y)$, give the actual equation
\[
                       F_H=\text{unit}\,[Y^2+X^2k(X,z)].
                                                            \tag{4.56c}
\]
The second $X$ follows from normal multiplicity two
along every nearby point of $A$.

\item The monic finite partial normalization is
\[
      N'=\{u^2+k(X,z)=0\},\qquad Y=Xu.
                                                            \tag{4.56d}
\]
It is an isomorphism away from $A$. The branch $F_H$
is generically smooth along $B_+$: the two generic
$B$ sheets belong separately to $F_H$ and $F_0$.
Thus $N'$ has at most an isolated singularity here.
Its hypersurface Cohen--Macaulay property and regularity
in codimension one make it normal, so it is the actual
normalization.

\item The reduced divisor $X=0$ is supported on $H$ and has
equation $u^2+k(0,z)=0$. The actual $H$ is smoothly
embedded and has degree two with ramification over $A$.
An even zero of $k(0,z)$ gives two branches; an odd
zero of order at least three gives a singular curve.
Therefore its zero is simple, $k_z(0,0)\ne0$, and
$N'$ is smooth. The minimal normalization resolution
is the identity in this neighbourhood.

\item Let $m\ge1$ be the generic intersection multiplicity of
the two branches along $B$. It is the same on both
local signs, because they belong to the same irreducible
image component. Let $G$ define $F_0$, and set
\[
                      b=\operatorname{ord}_{p,A}(G|_A)>0.
\]
On the smooth branch $F_R$, its divisor is $mB_-$.
Restricting to $A$ gives
\[
                 b=m I_{F_R}(A,B_-)\ge m.                \tag{4.56e}
\]
On the smooth normalization (4.56d), the divisor of $G$
is $mW$. The actual pole support is normal crossing,
so $H$ and $W$ meet transversely. Its restriction to
$H$ consequently has order $m$. But it is the pullback
of the same ambient restriction $G|_A$ under the local
degree-two map. Hence
\[
                              2b=m,
\]
contradicting (4.56e). Every case with $r_B\le2$ is
excluded.

\item The exact incidence partitions and the normalization
substitution are checked in
\texttt{code/part4/d4\_other\_two\_branches.py}.
Its corrected run reports PASS; the initial structural
symbolic-equality assertion failure is preserved separately.

\end{enumerate}
\end{proof}
